% Chapter 02 — eBPF Architecture
\section{Chương 2: Kiến trúc eBPF}
\ChapterDivider{Chương 2}{Kiến trúc eBPF}{Chương trình chạy ở đâu và các thành phần phối hợp thế nào?}

\begin{frame}{eBPF nằm ở đâu trong Linux?}
  \begin{block}{USER SPACE}
    Ứng dụng thông thường chạy ở đây. Công cụ như \texttt{bpftrace} hoặc
    ứng dụng sử dụng \texttt{libbpf} nạp chương trình và thu thập kết quả.
  \end{block}
  \vspace{-0.08cm}
  \begin{center}
    $\Big\downarrow\ \text{system call / BPF API}\ \Big\downarrow$
  \end{center}
  \vspace{-0.08cm}
  \begin{block}{KERNEL SPACE}
    Linux kernel quản lý tiến trình, bộ nhớ, tệp và mạng. Chương trình
    eBPF hoạt động tại những \textbf{hook được hỗ trợ} trong ngữ cảnh phù hợp.
  \end{block}
  \vspace{-0.08cm}
  \begin{center}
    $\Big\downarrow\ \text{quản lý thiết bị}\ \Big\downarrow$
  \end{center}
  \vspace{-0.08cm}
  \begin{block}{HARDWARE}
    CPU \quad | \quad RAM \quad | \quad Thiết bị I/O
  \end{block}
\end{frame}

\begin{frame}{Kiến trúc eBPF: hai giai đoạn khác nhau}
  \begin{block}{A. Chuẩn bị chương trình (trước khi xử lý sự kiện)}
    \centering
    \textbf{Userspace loader} $\longrightarrow$ \textbf{Load + Verifier}
    $\longrightarrow$ \textbf{eBPF Program} $\longrightarrow$ \textbf{Attach vào hook}
    \\[0.1cm]
    \footnotesize{JIT (nếu được hỗ trợ/bật) có thể dịch bytecode sang mã máy; không phải hook.}
  \end{block}
  \vspace{0.15cm}
  \begin{block}{B. Khi sự kiện xảy ra (thời điểm chạy)}
    \centering
    \textbf{Kernel event} $\longrightarrow$ \textbf{Hook}
    $\longrightarrow$ \textbf{eBPF Program} $\longrightarrow$ \textbf{Map / event output}
    \\[0.1cm]
    \textbf{Userspace application} $\longleftarrow$ \textbf{Kết quả}
  \end{block}
  \vspace{0.13cm}
  \begin{columns}[T,onlytextwidth]
    \begin{column}{0.49\textwidth}
      \begin{block}{Trong kernel}
        Program truy cập map qua helper hoặc cơ chế phù hợp.
      \end{block}
    \end{column}
    \begin{column}{0.49\textwidth}
      \begin{block}{Ở User Space}
        Ứng dụng truy cập map qua \texttt{bpf()} hoặc API như \texttt{libbpf}.
      \end{block}
    \end{column}
  \end{columns}
\end{frame}

\begin{frame}{Bốn thành phần quan trọng nhất}
  \begin{columns}[T,onlytextwidth]
    \begin{column}{0.49\textwidth}
      \begin{block}{eBPF Program}
        Chứa logic xử lý; chạy trong ngữ cảnh phù hợp với loại chương trình.
      \end{block}
      \begin{block}{Hook}
        Điểm gắn chương trình: ví dụ \textit{tracepoint}, \textit{kprobe}, XDP.
        Không phải program nào cũng gắn được vào mọi hook.
      \end{block}
    \end{column}
    \begin{column}{0.49\textwidth}
      \begin{block}{eBPF Map}
        Cấu trúc dữ liệu lưu trạng thái, bộ đếm hoặc chia sẻ thông tin
        với User Space.
      \end{block}
      \begin{block}{Verifier}
        Kiểm tra luồng điều khiển, truy cập bộ nhớ và thao tác được phép
        khi kernel nạp chương trình.
      \end{block}
    \end{column}
  \end{columns}
  \vspace{0.16cm}
  \begin{alertblock}{Một giới hạn cần nhớ}
    Verifier kiểm tra các điều kiện an toàn, \textbf{không} bảo đảm chương trình
    luôn đúng logic hoặc không gây chi phí xử lý.
  \end{alertblock}
\end{frame}

\begin{frame}{Điều gì xảy ra khi một sự kiện được ghi nhận?}
  \begin{block}{Ví dụ khái quát: một system call đi vào kernel}
    \centering
    \textbf{System call} $\longrightarrow$ \textbf{Tracepoint / Hook}
    $\longrightarrow$ \textbf{eBPF Program}\\[0.15cm]
    $\Big\downarrow$\\[0.12cm]
    \textbf{Lọc / lấy thông tin / cập nhật bộ đếm}\\[0.15cm]
    $\Big\downarrow$\\[0.12cm]
    \textbf{Map hoặc cơ chế gửi sự kiện} $\longrightarrow$ \textbf{User Space}
  \end{block}
  \vspace{0.2cm}
  \cue{1}{Chương trình phải được \textbf{nạp và gắn trước} khi sự kiện xảy ra.}
  \cue{2}{Một hook chỉ kích hoạt các chương trình được gắn phù hợp.}
  \cue{3}{Dữ liệu thu thập tùy vào program type, hook và logic được viết.}
\end{frame}
